\chapter{Computational Security}

In the previous chapter we always hit the same wall:
\begin{itemize}
  \item \textbf{perfect secrecy} $\to$ the key must be at least as long as the message (Shannon's theorem);
  \item \textbf{statistical MACs} $\to$ the key grows with the number of authenticated messages, $(t+1)\lambda$ bits;
  \item \textbf{randomness extraction} $\to$ from a source with min-entropy $k$ we can extract at most about $k$ bits.
\end{itemize}

To overcome these limits we accept \textbf{two assumptions}:
\begin{enumerate}
  \item the adversary is \textbf{computationally bounded}: it is an efficient (probabilistic polynomial-time) algorithm;
  \item there exist \textbf{hard problems}, that efficient adversaries cannot solve.
\end{enumerate}

Every result of this part of the course has the same shape: \textit{``if problem $X$ is hard for efficient algorithms, then scheme $\Pi$ is secure against efficient adversaries''}. It is proved by \textbf{reduction}: we show that anyone who breaks $\Pi$ efficiently could be turned into an efficient solver for $X$. The contrapositive is just as important: if $\Pi$ were ever broken, we would get an efficient algorithm for $X$, which for well-studied problems (factoring, discrete logarithm, \ldots) would be a breakthrough.

% ============================================================
%  One-way functions and Impagliazzo's worlds
% ============================================================

\section{One-way functions}

\subsection{The computational model}

From now on, ``efficient'' means \textbf{polynomial time} in the input length $n$: at most $n^c$ steps for some constant $c$, written $\poly(n)$. We are generous with the adversary: it can use $\poly(n)$ time and $\poly(n)$ \textbf{random bits}. In other words, the adversary is a \textbf{probabilistic polynomial-time} (PPT) algorithm. Its success probabilities are measured against \textbf{negligible functions} (\Cref{def:negl}): functions that go to zero faster than the inverse of any polynomial.

\subsection{The definition}

The hard problem underlying all of symmetric cryptography is the existence of functions that are \textbf{easy to compute but hard to invert}. Multiplying two large primes is immediate, while recovering the two primes from their product (factoring) is, as far as we know, extremely hard. It is like breaking an egg: easy in one direction, impossible in the other.

\begin{frameddefn}[label={def:owf}]{One-way function}
  A function $f : \{0,1\}^n \to \{0,1\}^m$ is a \textbf{one-way function} (OWF) if:
  \begin{enumerate}
    \item $f$ is computable in $\poly(n)$ time;
    \item for every PPT algorithm $\Adv$ there is a negligible function $\negl$ such that
    \[
      \Pr\big[f(x') = y \;:\; x \sample \{0,1\}^n,\ y = f(x),\ x' \leftarrow \Adv(y)\big] \leq \negl(n)
    \]
    where the probability is over the choice of $x$ and the random bits of $\Adv$.
  \end{enumerate}
\end{frameddefn}

How to read the definition:
\begin{itemize}
  \item \textbf{$x$ is chosen at random} $\to$ hardness is required \textbf{on average}, not only on a few particularly hard inputs. This is what cryptography needs, because keys are random.
  \item \textbf{Any preimage is enough} $\to$ the adversary wins if it finds \textit{some} $x'$ with $f(x') = y$, not necessarily the original $x$.
  \item \textbf{Negligible, not zero} $\to$ the adversary can always guess a random $x'$, so the success probability can never be exactly zero.
\end{itemize}

Nobody knows whether one-way functions exist: proving it would imply $\mathsf{P} \neq \mathsf{NP}$ (see below). There are, however, well-studied \textbf{candidates}:
\begin{itemize}
  \item \textbf{Multiplication} $\to$ $f(p, q) = p \cdot q$ for random primes $p, q$; inverting it means factoring.
  \item \textbf{Modular exponentiation} $\to$ $f(x) = g^x \bmod p$; inverting it means computing the discrete logarithm.
  \item \textbf{Subset sum} $\to$ $f(a_1, \ldots, a_k, S) = (a_1, \ldots, a_k, \sum_{i \in S} a_i)$; inverting it means finding a subset with a given sum.
\end{itemize}

\subsection{Why $\mathsf{P} \neq \mathsf{NP}$ is not enough}

One could ask: why not simply assume $\mathsf{P} \neq \mathsf{NP}$? Because it is too \textit{weak} for cryptography.
\begin{itemize}
  \item \textbf{OWFs imply $\mathsf{P} \neq \mathsf{NP}$.} Inverting $f$ is an $\mathsf{NP}$ search problem: given a candidate $x'$, checking that $f(x') = y$ is easy. If $\mathsf{P} = \mathsf{NP}$, every efficiently computable $f$ could be inverted, so no OWF would exist.
  \item \textbf{The converse is not known.} $\mathsf{P} \neq \mathsf{NP}$ only says that some problems are hard \textbf{in the worst case}, that is, on some inputs. Cryptography needs problems that are hard \textbf{on random inputs}, and with a way to generate hard instances together with their solutions. $\mathsf{P} \neq \mathsf{NP}$ does not guarantee any of this.
\end{itemize}

\subsection{Impagliazzo's five worlds}

In 1995 Russell Impagliazzo described five possible ``worlds'' we may be living in, depending on which complexity assumptions turn out to be true:

\begin{table}[H]
  \centering
  \small
  \begin{tabularx}{\textwidth}{@{}lX@{}}
    \toprule
    \textbf{World} & \textbf{What happens} \\
    \midrule
    Algorithmica & $\mathsf{P} = \mathsf{NP}$: every $\mathsf{NP}$ problem is easy, and no cryptography is possible \\
    Heuristica & $\mathsf{P} \neq \mathsf{NP}$, but hard problems are easy \textit{on average} \\
    Pessiland & there are problems that are hard on average, but no one-way functions: we pay the price of hardness without getting any cryptography \\
    Minicrypt & one-way functions exist (so symmetric cryptography is possible), but public-key cryptography is not \\
    Cryptomania & one-way functions \textit{and} public-key cryptography exist \\
    \bottomrule
  \end{tabularx}
  \caption{Impagliazzo's five worlds.}
\end{table}

Impagliazzo does not say which world we live in. Pessiland is the worst one, and most researchers believe (and hope) that we live in Cryptomania, or at least in Minicrypt.

% ============================================================
%  Computational indistinguishability and pseudorandom generators
% ============================================================

\section{Pseudorandom generators}

\subsection{Motivation: from a few random bits to many}

An extractor gives us a few truly random bits, but we often need many more: for example, a One-Time Pad key for a 1-gigabyte file. A \textbf{pseudorandom generator} (PRG) takes a short truly random string, the \textbf{seed}, and stretches it into a long string that \textbf{looks} random.

\textbf{Looks random to whom?} The output of a PRG \textbf{cannot} be truly uniform. If the seed has $n$ bits, the generator produces at most $2^n$ different strings out of the $2^{n+\ell}$ possible ones. An \textbf{unbounded} adversary can try all the seeds and check whether a given string is one of the possible outputs. So pseudorandomness makes sense \textbf{only against efficient adversaries}. This does not contradict Shannon's theorem, which holds against unbounded adversaries.

\subsection{Computational indistinguishability}

We need the computational counterpart of the statistical distance (\Cref{def:sd}). Since security is asymptotic in $n$, we talk about \textbf{probability ensembles}: sequences $\mathcal{X} = \{X_n\}_{n \in \N}$ of distributions, one for every value of the security parameter.

\begin{frameddefn}[label={def:cind}]{Computational indistinguishability}
  Two ensembles $\mathcal{X} = \{X_n\}_{n \in \N}$ and $\mathcal{Y} = \{Y_n\}_{n \in \N}$ are \textbf{computationally indistinguishable}, written $\mathcal{X} \cind \mathcal{Y}$, if for every PPT algorithm $D$ (the \textit{distinguisher}) there is a negligible function $\negl$ such that, for every $n$:
  \[
    \Big|\Pr[D(z) = 1 : z \leftarrow X_n] - \Pr[D(z) = 1 : z \leftarrow Y_n]\Big| \leq \negl(n)
  \]
\end{frameddefn}

It is the same ``quiz'' as always: $D$ receives a string and must say which distribution it comes from, and it must not do noticeably better than tossing a coin. The differences with the statistical distance are that the distinguisher must be \textbf{efficient}, and that the advantage must be \textbf{negligible} as $n$ grows.

Useful properties:
\begin{itemize}
  \item \textbf{Efficient processing cannot help} $\to$ if $\mathcal{X} \cind \mathcal{Y}$ and $f$ is computable in polynomial time, then $f(\mathcal{X}) \cind f(\mathcal{Y})$. Otherwise, a distinguisher for $f(\mathcal{X})$ and $f(\mathcal{Y})$ composed with $f$ would distinguish $\mathcal{X}$ from $\mathcal{Y}$.
  \item \textbf{Transitivity} $\to$ if $\mathcal{X} \cind \mathcal{Y}$ and $\mathcal{Y} \cind \mathcal{Z}$, then $\mathcal{X} \cind \mathcal{Z}$, because the sum of two negligible functions is negligible. This extends to a \textit{constant} number of steps; for a polynomial number of steps we need the hybrid argument of \Cref{thm:stretch}.
  \item \textbf{Statistical closeness implies computational indistinguishability} $\to$ if $\SD(X_n; Y_n)$ is negligible, then $\mathcal{X} \cind \mathcal{Y}$, since the statistical distance bounds the advantage of every distinguisher, efficient or not.
\end{itemize}

\subsection{The definition of PRG}

\begin{frameddefn}[label={def:prg}]{Pseudorandom generator}
  A function $G : \{0,1\}^n \to \{0,1\}^{n+\ell}$, with $\ell \geq 1$, is a \textbf{pseudorandom generator} with stretch $\ell$ if:
  \begin{enumerate}
    \item $G$ is computable in $\poly(n)$ time;
    \item $G(\Uni{n}) \cind \Uni{n+\ell}$, where $\Uni{k}$ is the uniform distribution over $\{0,1\}^k$.
  \end{enumerate}
\end{frameddefn}

Observations:
\begin{itemize}
  \item \textbf{$G$ is deterministic} $\to$ all the randomness of the output comes from the seed. The strength of the definition is that $n$ random bits become $n + \ell$ bits that no efficient algorithm can tell apart from uniform.
  \item \textbf{The seed must be uniform} (or at least computationally indistinguishable from uniform). In practice, a randomness extractor produces the seed, and the PRG expands it.
  \item \textbf{A first application.} With a PRG, Alice and Bob can share a short seed $s$ and use $G(s)$ as the pad of a One-Time Pad for a much longer message. The scheme is not perfectly secret (Shannon's theorem forbids it), but it is secure against efficient adversaries. We will make this precise in the next chapters.
\end{itemize}

\subsection{Stretching a PRG}

How much should a PRG stretch? Surprisingly, it is enough to build a PRG that adds \textbf{a single bit}: from it we can obtain a PRG that adds any polynomial number of bits.

\paragraph{The construction.} Let $G : \{0,1\}^n \to \{0,1\}^{n+1}$ be a PRG, and split its output as $G(s) = (s', b)$ with $s' \in \{0,1\}^n$ and $b \in \{0,1\}$. Apply $G$ to the seed, keep the last bit as output, and apply $G$ again to the remaining $n$ bits:
\[
  s_0 \xrightarrow{\ G\ } (s_1, b_1), \qquad s_1 \xrightarrow{\ G\ } (s_2, b_2), \qquad \ldots, \qquad s_{\ell-1} \xrightarrow{\ G\ } (s_\ell, b_\ell)
\]
and define
\[
  G^\ell(s_0) = b_1 \concat b_2 \concat \cdots \concat b_\ell \concat s_\ell \in \{0,1\}^{n+\ell}
\]

\begin{figure}[H]
  \centering
  \begin{tikzpicture}[g/.style={draw, rounded corners=2pt, minimum size=6mm}, x=1.15cm]
    \node (s0) at (0,0) {$s_0$};
    \node[g] (g1) at (1.2,0) {$G$};
    \node (s1) at (2.4,0) {$s_1$};
    \node[g] (g2) at (3.6,0) {$G$};
    \node (s2) at (4.8,0) {$s_2$};
    \node (dots) at (6,0) {$\cdots$};
    \node[g] (gl) at (7.2,0) {$G$};
    \node (sl) at (8.4,0) {$s_\ell$};
    \node (b1) at (1.2,-1.2) {$b_1$};
    \node (b2) at (3.6,-1.2) {$b_2$};
    \node (bl) at (7.2,-1.2) {$b_\ell$};
    \draw[msg] (s0) -- (g1);
    \draw[msg] (g1) -- (s1);
    \draw[msg] (s1) -- (g2);
    \draw[msg] (g2) -- (s2);
    \draw[msg] (s2) -- (dots);
    \draw[msg] (dots) -- (gl);
    \draw[msg] (gl) -- (sl);
    \draw[msg] (g1) -- (b1);
    \draw[msg] (g2) -- (b2);
    \draw[msg] (gl) -- (bl);
  \end{tikzpicture}
  \caption{The stretching construction: every application of $G$ outputs one bit and a new state.}
\end{figure}

\begin{framedthm}[label={thm:stretch}]{PRG stretching}
  If $G : \{0,1\}^n \to \{0,1\}^{n+1}$ is a PRG, then for every $\ell = \poly(n)$ the function $G^\ell : \{0,1\}^n \to \{0,1\}^{n+\ell}$ defined above is a PRG.
\end{framedthm}

The proof uses the \textbf{hybrid argument}, one of the most important techniques of the course. The idea is to build a chain of distributions that goes, one small step at a time, from the output of $G^\ell$ to a truly uniform string.

\paragraph{The hybrids.} For $i = 0, \ldots, \ell$, the hybrid $H_i$ is the distribution of the output of this procedure:
\begin{enumerate}
  \item choose the first $i$ output bits $b_1, \ldots, b_i$ uniformly at random, and a uniform state $s_i \sample \{0,1\}^n$;
  \item for $j = i+1, \ldots, \ell$, compute $(s_j, b_j) = G(s_{j-1})$;
  \item output $b_1 \concat \cdots \concat b_\ell \concat s_\ell$.
\end{enumerate}
In words, in $H_i$ the first $i$ bits are \textbf{truly random}, and from there on the string is produced by $G$.

\begin{table}[H]
  \centering
  \small
  \begin{tabular}{lccccc}
    \toprule
    \textbf{Hybrid} & $b_1$ & $b_2$ & $\cdots$ & $b_\ell$ & $s_\ell$ \\
    \midrule
    $H_0 = G^\ell(\Uni{n})$ & $G$ & $G$ & $\cdots$ & $G$ & $G$ \\
    $H_1$ & random & $G$ & $\cdots$ & $G$ & $G$ \\
    $H_2$ & random & random & $\cdots$ & $G$ & $G$ \\
    $\vdots$ & & & & & \\
    $H_\ell = \Uni{n+\ell}$ & random & random & $\cdots$ & random & random \\
    \bottomrule
  \end{tabular}
  \caption{The hybrids: ``$G$'' means produced by the chain of $G$, ``random'' means truly uniform. In $H_\ell$ the final state $s_\ell$ is the uniform state chosen in step 1.}
\end{table}

The two extremes are exactly the distributions we want to compare: $H_0 = G^\ell(\Uni{n})$ and $H_\ell = \Uni{n+\ell}$.

\begin{proof}
  $G^\ell$ is computable in polynomial time, since it runs $G$ exactly $\ell = \poly(n)$ times. We must show that $H_0 \cind H_\ell$.

  \textbf{Step 1: neighbouring hybrids differ in one application of $G$.} Compare $H_i$ and $H_{i+1}$. In both, $b_1, \ldots, b_i$ are random and $s_i$ is uniform. The only difference is the pair $(s_{i+1}, b_{i+1})$:
  \begin{itemize}
    \item in $H_i$ it is $G(s_i)$ with $s_i$ uniform, that is, a sample of $G(\Uni{n})$;
    \item in $H_{i+1}$ it is a uniform string of $n+1$ bits, that is, a sample of $\Uni{n+1}$.
  \end{itemize}
  Everything after that pair is computed in the same way, by applying $G$.

  \textbf{Step 2: a distinguisher for neighbours becomes a distinguisher for $G$.} Let $D$ be any PPT algorithm. For each $i$, define the algorithm $D_i$ that, on input $w = (s', b') \in \{0,1\}^{n+1}$:
  \begin{enumerate}
    \item chooses $b_1, \ldots, b_i$ uniformly at random;
    \item sets $s_{i+1} = s'$ and $b_{i+1} = b'$;
    \item computes $(s_j, b_j) = G(s_{j-1})$ for $j = i+2, \ldots, \ell$;
    \item returns $D(b_1 \concat \cdots \concat b_\ell \concat s_\ell)$.
  \end{enumerate}
  If $w \leftarrow G(\Uni{n})$, the string given to $D$ is distributed exactly as $H_i$; if $w \leftarrow \Uni{n+1}$, it is distributed exactly as $H_{i+1}$. Writing $p_i = \Pr[D(z) = 1 : z \leftarrow H_i]$:
  \[
    \Pr[D_i(w) = 1 : w \leftarrow G(\Uni{n})] - \Pr[D_i(w) = 1 : w \leftarrow \Uni{n+1}] = p_i - p_{i+1}
  \]

  \textbf{Step 3: from neighbours to the extremes.} Let $D^*$ be the algorithm that chooses $i \in \{0, \ldots, \ell-1\}$ uniformly at random and runs $D_i$. It is PPT because $\ell$ is polynomial. Its advantage against $G$ is the average of the $\ell$ differences, and the sum \textbf{telescopes}:
  \[
    \Pr[D^*(w) = 1 : w \leftarrow G(\Uni{n})] - \Pr[D^*(w) = 1 : w \leftarrow \Uni{n+1}]
    = \frac{1}{\ell}\sum_{i=0}^{\ell-1} (p_i - p_{i+1}) = \frac{p_0 - p_\ell}{\ell}
  \]
  Since $G$ is a PRG, the left side is at most $\negl(n)$ in absolute value. Therefore
  \[
    \big|\Pr[D(z) = 1 : z \leftarrow G^\ell(\Uni{n})] - \Pr[D(z) = 1 : z \leftarrow \Uni{n+\ell}]\big| = |p_0 - p_\ell| \leq \ell \cdot \negl(n)
  \]
  and $\ell \cdot \negl(n)$ is still negligible because $\ell$ is a polynomial. Since $D$ was arbitrary, $G^\ell(\Uni{n}) \cind \Uni{n+\ell}$.
\end{proof}

\begin{framedobs}{Why $\ell$ must be polynomial}
  The advantage between the two extremes is at most the \textbf{sum} of the advantages between neighbours. With $\ell$ hybrids this is at most $\ell \cdot \negl(n)$, which is negligible only if $\ell$ is polynomial: a polynomial times a negligible function is negligible, while an exponential times a negligible function may not be. This is why ``transitivity of $\cind$'' alone is not enough here: it would only cover a constant number of steps.
\end{framedobs}

% ============================================================
%  From PRGs to OWFs, hard-core predicates, from OWFs to PRGs
% ============================================================

\section{Hard-core predicates}

\Cref{thm:stretch} tells us that a PRG adding a single bit is enough. How do we build one? We will see that PRGs and one-way functions are \textbf{equivalent}: one exists if and only if the other does. Since nobody knows whether one-way functions exist, building a PRG is exactly as hard as building a OWF; in practice, real-world PRGs are \textit{heuristic} constructions, believed but not proved to be secure.

\subsection{From PRGs to one-way functions}

The easy direction first.

\begin{framedthm}[label={thm:prg-owf}]{From PRGs to OWFs}
  Let $G : \{0,1\}^n \to \{0,1\}^{2n}$ be a PRG. Then the function $f : \{0,1\}^{2n} \to \{0,1\}^{2n}$ defined by
  \[
    f(x_1 \concat x_2) = G(x_1) \qquad \text{with } x_1, x_2 \in \{0,1\}^n
  \]
  is a one-way function.
\end{framedthm}

\textbf{Intuition.} Whoever can invert $f$ can \textbf{recognize} the outputs of $G$: take the string, try to invert it, and recompute. On outputs of $G$ the inversion works with noticeable probability, while a uniform string is almost never an output of $G$ at all.

\begin{proof}
  $f$ is efficiently computable because $G$ is. Suppose, by contradiction, that a PPT adversary $\Adv$ inverts $f$ with a probability $\varepsilon(n)$ that is not negligible. Build the distinguisher $D$ that, on input $z \in \{0,1\}^{2n}$:
  \begin{enumerate}
    \item computes $x' \leftarrow \Adv(z)$;
    \item outputs $1$ if $f(x') = z$, and $0$ otherwise.
  \end{enumerate}
  \begin{itemize}
    \item \textbf{If $z \leftarrow G(\Uni{n})$}: since $f$ ignores the second half of its input, $z = G(x_1)$ with $x_1$ uniform is distributed exactly as $f(x)$ with $x$ uniform. So $D$ outputs $1$ exactly when $\Adv$ inverts, that is, with probability $\varepsilon(n)$.
    \item \textbf{If $z \leftarrow \Uni{2n}$}: $D$ outputs $1$ only if $z$ is in the image of $G$. The image contains at most $2^n$ strings out of $2^{2n}$, so this happens with probability at most $2^{-n}$.
  \end{itemize}
  The advantage of $D$ is at least $\varepsilon(n) - 2^{-n}$, which is not negligible (for large $n$ it is at least $\frac{\varepsilon(n)}{2}$). This contradicts the fact that $G$ is a PRG.
\end{proof}

\subsection{Hard-core predicates}

For the opposite direction we need a bit of the input that stays \textbf{secret} even when $f(x)$ is known.

\begin{frameddefn}[label={def:hardcore}]{Hard-core predicate}
  Let $f : \{0,1\}^n \to \{0,1\}^m$ be efficiently computable. A function $h : \{0,1\}^n \to \{0,1\}$ is a \textbf{hard-core predicate} for $f$ if for every PPT adversary $\Adv$:
  \[
    \Pr\big[\Adv(f(x)) = h(x) : x \sample \{0,1\}^n\big] \leq \frac{1}{2} + \negl(n)
  \]
\end{frameddefn}

In words: seeing $f(x)$, no efficient adversary can guess the bit $h(x)$ noticeably better than tossing a coin. An equivalent formulation, which is often more convenient, is
\[
  \big(f(x),\ h(x)\big) \cind \big(f(x),\ \Uni{1}\big) \qquad \text{with } x \sample \{0,1\}^n
\]
that is, $h(x)$ looks like a fresh random bit even next to $f(x)$.

Neither definition requires $f$ to be one-way, but intuitively a hard-core bit is what makes a function hard to invert. This raises three natural questions.

\paragraph{Question 1: is there a universal hard-core predicate, hard-core for every OWF?} \textbf{No.}

\begin{proof}
  Suppose that an efficiently computable $h$ were hard-core for every OWF. Take any OWF $f$ and build a ``bad'' function that \textbf{publishes} the bit $h(x)$ next to $f(x)$:
  \[
    f_{\mathrm{BAD}}(x) = f(x) \concat h(x)
  \]
  \begin{itemize}
    \item \textbf{$h$ is clearly not hard-core for $f_{\mathrm{BAD}}$}: the adversary that outputs the last bit of $f_{\mathrm{BAD}}(x)$ guesses $h(x)$ with probability $1$.
    \item \textbf{$f_{\mathrm{BAD}}$ is still one-way.} We prove it by reduction. Suppose that there is a PPT adversary $\Adv$ that inverts $f_{\mathrm{BAD}}$ with probability at least $\frac{1}{\poly(n)}$, that is, not negligible. We build a PPT adversary $B$ that inverts $f$, playing between the challenger $C$ of the one-wayness game of $f$ and the adversary $\Adv$:
  \end{itemize}

  \begin{center}
    \begin{tikzpicture}[x=1cm, y=1cm]
      \node (A) at (0,0) {$\Adv$};
      \node (B) at (4.5,0) {$B$};
      \node (C) at (9,0) {$C$};
      \node[lbl, align=left] at (9,-1.2) {$x \sample \{0,1\}^n$};
      \draw[msg] (9,-0.6) -- node[above] {$y = f(x)$} (4.5,-0.6);
      \node[lbl] at (4.5,-1.2) {$\tilde b \sample \{0,1\}$};
      \draw[msg] (4.5,-1.6) -- node[above] {$y \concat \tilde b$} (0,-1.6);
      \draw[msg] (0,-2.4) -- node[above] {$x'$} (4.5,-2.4);
      \draw[msg] (4.5,-2.4) -- node[above] {$x'$} (9,-2.4);
    \end{tikzpicture}
  \end{center}

  $B$ receives $y = f(x)$, \textbf{guesses} the missing bit with a coin $\tilde b$, gives $y \concat \tilde b$ to $\Adv$ and forwards its answer $x'$.
  \begin{itemize}
    \item With probability $\frac{1}{2}$ the guess is right, $\tilde b = h(x)$. The coin is independent of $x$, so in this case the input of $\Adv$ is distributed \textbf{exactly} as $f_{\mathrm{BAD}}(x)$ for a random $x$, and $\Adv$ finds $x'$ with $f_{\mathrm{BAD}}(x') = y \concat \tilde b$ with probability at least $\frac{1}{\poly(n)}$.
    \item $f_{\mathrm{BAD}}(x') = y \concat \tilde b$ implies in particular $f(x') = y$, so $B$ wins.
  \end{itemize}
  Therefore
  \[
    \Pr[B \text{ wins}] \geq \frac{1}{2} \cdot \Pr[\Adv \text{ wins}] \geq \frac{1}{2} \cdot \frac{1}{\poly(n)} = \frac{1}{\poly'(n)}
  \]
  which is not negligible, against the one-wayness of $f$. (The inequality is $\geq$ and not $=$ because $B$ may win even when the guess is wrong.)

  So $f_{\mathrm{BAD}}$ is a OWF for which $h$ is not hard-core, a contradiction.
\end{proof}

\textbf{The next best thing.} Since no single $h$ works for every $f$, we reverse the quantifiers: instead of ``$\exists\, h\ \forall f$'' we look for ``$\forall f\ \exists\, h$'', that is, a hard-core predicate tailored to each function. Question 3 shows that this is possible, up to a small modification of $f$.

\paragraph{Question 2: is every function with a hard-core predicate one-way?} \textbf{No.} Let $f$ delete the last bit, $f(x_1 \cdots x_n) = x_1 \cdots x_{n-1}$, and let $h(x) = x_n$ be the deleted bit. Since $x$ is uniform, $x_n$ is independent of $f(x)$, so $h$ is hard-core (no adversary can do better than $\frac{1}{2}$). But $f$ is trivially invertible: appending any bit to $y$ gives a valid preimage, so the adversary inverts with probability $1$.

\paragraph{Question 3: does every OWF have a hard-core predicate?} \textbf{Yes, up to a small modification.} This is the Goldreich--Levin theorem.

\begin{framedthm}[label={thm:gl}]{Goldreich--Levin}
  Let $f : \{0,1\}^n \to \{0,1\}^m$ be a OWF, and define $g : \{0,1\}^{2n} \to \{0,1\}^{m+n}$ as
  \[
    g(x, r) = f(x) \concat r
  \]
  Then $g$ is a OWF, and it has the hard-core predicate
  \[
    h(x, r) = \langle x, r \rangle \bmod 2 = x_1 r_1 \oplus x_2 r_2 \oplus \cdots \oplus x_n r_n
  \]
\end{framedthm}

How to read it:
\begin{itemize}
  \item \textbf{$r$ is public} $\to$ it appears in clear in the output $g(x, r)$, exactly like the seed of an extractor.
  \item \textbf{$h(x, r)$ is the XOR of a random subset of the bits of $x$} $\to$ the subset is selected by $r$: bit $x_i$ is included if and only if $r_i = 1$. It is the same family of functions as in the Alice and Bob extraction example.
  \item \textbf{$g$ is one-way} $\to$ the second half $r$ is in clear and carries no information about $x$, so inverting $g$ means inverting $f$.
\end{itemize}

\paragraph{Intuition: a perfect predictor would invert $f$.} Suppose there were a very good adversary $\Adv$ that, given $g(x, r) = f(x) \concat r$, always computes $h(x, r) = \langle x, r \rangle$, \textbf{for every} $r$. Then we could recover $x$ bit by bit, by choosing as $r$ the unit vectors $e_i = (0, \ldots, 0, 1, 0, \ldots, 0)$, with the $1$ in position $i$:
\[
  \langle x, e_1 \rangle = \langle x, 100\cdots0 \rangle = x_1, \qquad \ldots, \qquad \langle x, e_n \rangle = \langle x, 0\cdots01 \rangle = x_n
\]
So $n$ calls $\Adv(f(x) \concat e_i)$ would return all the bits of $x$, and $f$ would be inverted, against the hypothesis.

\paragraph{Why the real proof is harder.} The hypothesis only gives an adversary that guesses $h(x, r)$ with probability $\frac{1}{2} + \frac{1}{\poly(n)}$, and only \textbf{on average over a random $r$}. Such an adversary may be wrong on every unit vector $e_i$, which are very particular values of $r$. The full proof (omitted in class) has to work with random $r$ only.

\begin{framedobs}{Beyond the lecture: the self-correction trick}
  An intermediate case shows how to avoid the particular values $e_i$. Suppose that $\Adv$ guesses $\langle x, r \rangle$ with probability at least $\frac{3}{4} + \varepsilon$ over a random $r$. Since the inner product is linear,
  \[
    \langle x, r \rangle \oplus \langle x, r \oplus e_i \rangle = \langle x, e_i \rangle = x_i
  \]
  So, to compute $x_i$, choose a \textbf{random} $r$ and ask $\Adv$ both $\langle x, r \rangle$ and $\langle x, r \oplus e_i \rangle$. Both $r$ and $r \oplus e_i$ are uniformly distributed, so each answer is wrong with probability at most $\frac{1}{4} - \varepsilon$, and by the union bound both are right with probability at least $\frac{1}{2} + 2\varepsilon$. Repeating with many independent $r$ and taking the majority gives $x_i$ with high probability. The real proof, for probability $\frac{1}{2} + \varepsilon$, needs a cleverer way of choosing the queries.
\end{framedobs}

\subsection{From one-way permutations to PRGs}

Hard-core predicates are exactly what we need to build a PRG. The simplest case uses a \textbf{one-way permutation} (OWP): a OWF that is a bijection $f : \{0,1\}^n \to \{0,1\}^n$.

\begin{framedprop}[label={prop:owp-prg}]{From OWPs to PRGs}
  Let $f$ be a one-way permutation with hard-core predicate $h$. Then
  \[
    G(s) = f(s) \concat h(s)
  \]
  is a PRG with stretch $1$.
\end{framedprop}

\begin{proof}
  $G$ is efficiently computable because $f$ and $h$ are. For $s$ uniform:
  \begin{itemize}
    \item $f(s)$ is \textbf{uniform} over $\{0,1\}^n$, because a permutation of a uniform string is uniform;
    \item by the equivalent formulation of hard-core predicates, $h(s)$ looks like a fresh random bit even next to $f(s)$.
  \end{itemize}
  Therefore
  \[
    G(\Uni{n}) = \big(f(s),\, h(s)\big) \cind \big(f(s),\, \Uni{1}\big) = \big(\Uni{n},\, \Uni{1}\big) = \Uni{n+1} \qedhere
  \]
\end{proof}

Combined with Goldreich--Levin: if $f$ is a OWP, then $g(x, r) = f(x) \concat r$ is also a permutation, so $G(x, r) = f(x) \concat r \concat \langle x, r \rangle$ is a PRG with stretch $1$.

\subsection{From one-way functions to PRGs}

The proposition can be extended to \textbf{every} one-way function, not only permutations. The construction is much harder (it combines many evaluations of $f$ with their hard-core bits), and its proof was omitted.

\begin{framedthm}[label={thm:hill}]{Håstad, Impagliazzo, Levin, Luby}
  If a one-way function exists, then a PRG with stretch $1$ exists.
\end{framedthm}

Together with \Cref{thm:prg-owf} and \Cref{thm:stretch} this gives the complete picture:
\[
  \begin{aligned}
    \text{OWF} \;&\xrightarrow{\ \text{Goldreich--Levin, HILL}\ }\; \text{PRG with stretch } 1 \;\xrightarrow{\ \text{stretching}\ }\; \text{PRG with stretch } \poly(n) \\
    \text{PRG} \;&\xrightarrow{\ \text{\Cref{thm:prg-owf}}\ }\; \text{OWF}
  \end{aligned}
\]
\begin{framedcor}{}
  One-way functions exist if and only if pseudorandom generators exist.
\end{framedcor}

